Tabulka~\ref{tab:flexicurity} shrnuje šestici srovnávaných ekonomik. \Ac{ČR} dosahuje v~souboru nejvyšší míry zaměstnanosti (\SI{82,9}{\percent}) a~spolu se~Slovenskem nejnižšího Giniho koeficientu (\SI{24,0}{\percent}); současně však vykazuje nejdelší týdenní pracovní dobu (\SI{39,2}{\hour\per\week}) a~úplné náklady práce \SI{24,5}{\pps\per\hour}, tj.~\SI{61}{\percent} dánské či německé hodnoty. Daňový klín \aclgen{ČR} při \SI{67}{\percent} průměrné mzdy dosahuje \SI{38,9}{\percent} oproti \SI{33,5}{\percent} v~Dánsku.

Disponibilní hodinový příjem zaměstnance v~\aclloc{PPS} dosahuje \SI{176}{\percent} české úrovně v~Dánsku, \SI{155}{\percent} v~Rakousku, \SI{151}{\percent} v~Německu, zatímco Polsko (\SI{108}{\percent}) a~Slovensko (\SI{101}{\percent}) zůstávají v~pásmu \aclgen{ČR}. Podíl nízkopříjmových zaměstnanců (méně než~\SI{67}{\percent} mediánu) činí v~\acloc{ČR} \SI{13,8}{\percent} proti \SI{19,0}{\percent} v~Polsku i~Německu a~\SI{9,7}{\percent} v~Dánsku; volná pracovní místa činí \SI{1,6}{\percent} proti \SI{2,2}{\percent} v~Německu.

Trojice s~vysokým pokrytím \acgen{KS} (Rakousko \SI{98}{\percent}, Dánsko \SI{82}{\percent}, Německo \SI{54}{\percent}) vykazuje kratší pracovní dobu i~vyšší výdaje na \ac{APZ}, zatímco \ac{ČR}, Polsko a~Slovensko zůstávají v~pásmu pokrytí \SIrange{12}{31}{\percent}, hustoty odborů \SIrange{9}{12}{\percent} a~výdajů na \acs{APZ} pod \SI{0,21}{\percent}~\acs{HDP}.

Tabulka~\ref{tab:flexicurity} shrnuje šestici srovnávaných ekonomik. \Ac{ČR} dosahuje v~souboru nejvyšší míry zaměstnanosti (\SI{82,9}{\percent}) a~spolu se~Slovenskem nejnižšího Giniho koeficientu (\SI{24,0}{\percent}); současně však vykazuje nejdelší týdenní pracovní dobu (\SI{39,2}{\hour\per\week}) a~úplné náklady práce \SI{24,5}{\pps\per\hour}, tj.~\SI{61}{\percent} dánské či německé hodnoty. Daňový klín \aclgen{ČR} při \SI{67}{\percent} průměrné mzdy dosahuje \SI{38,9}{\percent} oproti \SI{33,5}{\percent} v~Dánsku.

Disponibilní hodinový příjem zaměstnance v~\aclloc{PPS} dosahuje \SI{176}{\percent} české úrovně v~Dánsku, \SI{155}{\percent} v~Rakousku, \SI{151}{\percent} v~Německu, zatímco Polsko (\SI{108}{\percent}) a~Slovensko (\SI{101}{\percent}) zůstávají v~pásmu \aclgen{ČR}. Podíl nízkopříjmových zaměstnanců (méně než~\SI{67}{\percent} mediánu) činí v~\acloc{ČR} \SI{13,8}{\percent} proti \SI{19,0}{\percent} v~Polsku i~Německu a~\SI{9,7}{\percent} v~Dánsku; volná pracovní místa činí \SI{1,6}{\percent} proti \SI{2,2}{\percent} v~Německu.

Trojice s~vysokým pokrytím \acgen{KS} (Rakousko \SI{98}{\percent}, Dánsko \SI{82}{\percent}, Německo \SI{54}{\percent}) vykazuje kratší pracovní dobu i~vyšší výdaje na \ac{APZ}, zatímco \ac{ČR}, Polsko a~Slovensko zůstávají v~pásmu pokrytí \SIrange{12}{31}{\percent}, hustoty odborů \SIrange{9}{12}{\percent} a~výdajů na \acs{APZ} pod \SI{0,21}{\percent}~\acs{HDP}.

Tabulka~\ref{tab:flexicurity} shrnuje šestici srovnávaných ekonomik. \Ac{ČR} dosahuje v~souboru nejvyšší míry zaměstnanosti (\SI{82,9}{\percent}) a~spolu se~Slovenskem nejnižšího Giniho koeficientu (\SI{24,0}{\percent}); současně však vykazuje nejdelší týdenní pracovní dobu (\SI{39,2}{\hour\per\week}) a~úplné náklady práce \SI{24,5}{\pps\per\hour}, tj.~\SI{61}{\percent} dánské či německé hodnoty. Daňový klín \aclgen{ČR} při \SI{67}{\percent} průměrné mzdy dosahuje \SI{38,9}{\percent} oproti \SI{33,5}{\percent} v~Dánsku.

Disponibilní hodinový příjem zaměstnance v~\aclloc{PPS} dosahuje \SI{176}{\percent} české úrovně v~Dánsku, \SI{155}{\percent} v~Rakousku, \SI{151}{\percent} v~Německu, zatímco Polsko (\SI{108}{\percent}) a~Slovensko (\SI{101}{\percent}) zůstávají v~pásmu \aclgen{ČR}. Podíl nízkopříjmových zaměstnanců (méně než~\SI{67}{\percent} mediánu) činí v~\acloc{ČR} \SI{13,8}{\percent} proti \SI{19,0}{\percent} v~Polsku i~Německu a~\SI{9,7}{\percent} v~Dánsku; volná pracovní místa činí \SI{1,6}{\percent} proti \SI{2,2}{\percent} v~Německu.

Trojice s~vysokým pokrytím \acgen{KS} (Rakousko \SI{98}{\percent}, Dánsko \SI{82}{\percent}, Německo \SI{54}{\percent}) vykazuje kratší pracovní dobu i~vyšší výdaje na \ac{APZ}, zatímco \ac{ČR}, Polsko a~Slovensko zůstávají v~pásmu pokrytí \SIrange{12}{31}{\percent}, hustoty odborů \SIrange{9}{12}{\percent} a~výdajů na \acs{APZ} pod \SI{0,21}{\percent}~\acs{HDP}.

Tabulka~\ref{tab:flexicurity} shrnuje šestici srovnávaných ekonomik. \Ac{ČR} dosahuje v~souboru nejvyšší míry zaměstnanosti (\SI{82,9}{\percent}) a~spolu se~Slovenskem nejnižšího Giniho koeficientu (\SI{24,0}{\percent}); současně však vykazuje nejdelší týdenní pracovní dobu (\SI{39,2}{\hour\per\week}) a~úplné náklady práce \SI{24,5}{\pps\per\hour}, tj.~\SI{61}{\percent} dánské či německé hodnoty. Daňový klín \aclgen{ČR} při \SI{67}{\percent} průměrné mzdy dosahuje \SI{38,9}{\percent} oproti \SI{33,5}{\percent} v~Dánsku.

Disponibilní hodinový příjem zaměstnance v~\aclloc{PPS} dosahuje \SI{176}{\percent} české úrovně v~Dánsku, \SI{155}{\percent} v~Rakousku, \SI{151}{\percent} v~Německu, zatímco Polsko (\SI{108}{\percent}) a~Slovensko (\SI{101}{\percent}) zůstávají v~pásmu \aclgen{ČR}. Podíl nízkopříjmových zaměstnanců (méně než~\SI{67}{\percent} mediánu) činí v~\acloc{ČR} \SI{13,8}{\percent} proti \SI{19,0}{\percent} v~Polsku i~Německu a~\SI{9,7}{\percent} v~Dánsku; volná pracovní místa činí \SI{1,6}{\percent} proti \SI{2,2}{\percent} v~Německu.

Trojice s~vysokým pokrytím \acgen{KS} (Rakousko \SI{98}{\percent}, Dánsko \SI{82}{\percent}, Německo \SI{54}{\percent}) vykazuje kratší pracovní dobu i~vyšší výdaje na \ac{APZ}, zatímco \ac{ČR}, Polsko a~Slovensko zůstávají v~pásmu pokrytí \SIrange{12}{31}{\percent}, hustoty odborů \SIrange{9}{12}{\percent} a~výdajů na \acs{APZ} pod \SI{0,21}{\percent}~\acs{HDP}.

\input{texparts/python/prakticka_srovnani}



